%  LaTeX support: latex@mdpi.com 
%  For support, please attach all files needed for compiling as well as the log file, and specify your operating system, LaTeX version, and LaTeX editor.

%=================================================================
\documentclass[computers,article,submit,pdftex,moreauthors]{Definitions/mdpi} 
%\documentclass[preprints,article,submit,pdftex,moreauthors]{Definitions/mdpi} 
% For posting an early version of this manuscript as a preprint, you may use "preprints" as the journal. Changing "submit" to "accept" before posting will remove line numbers.

% Below journals will use APA reference format:
% admsci, aieduc, behavsci, businesses, econometrics, economies, education, ejihpe, famsci, games, humans, ijcs, ijfs, jintelligence, journalmedia, jrfm, jsam, languages, peacestud, psycholint, publications, tourismhosp, youth

% Below journals will use Chicago reference format:
% arts, genealogy, histories, humanities, laws, literature, religions, risks, socsci

%--------------------
% Class Options:
%--------------------
%----------
% journal
%----------
% Choose between the following MDPI journals:
% accountaudit, acoustics, actuators, addictions, adhesives, admsci, adolescents, aerobiology, aerospace, agriculture, agriengineering, agrochemicals, agronomy, ai, aichem, aieduc, aieng, aimater, aimed, aipa, air, aisens, algorithms, allergies, alloys, amh, analog, analytica, analytics, anatomia, anesthres, animals, antibiotics, antibodies, antioxidants, applbiosci, appliedchem, appliedmath, appliedphys, applmech, applmicrobiol, applnano, applsci, aquacj, architecture, arm, arthropoda, arts, asc, asi, astronautics, astronomy, atmosphere, atoms, audiolres, automation, axioms, bacteria, batteries, bdcc, behavsci, beverages, biochem, bioengineering, biologics, biology, biomass, biomechanics, biomed, biomedicines, biomedinformatics, biomimetics, biomolecules, biophysica, bioresourbioprod, biosensors, biosphere, biotech, birds, blockchains, bloods, blsf, brainsci, breath, buildings, businesses, cancers, carbon, cardio, cardiogenetics, cardiovascmed, catalysts, cells, ceramics, challenges, chemengineering, chemistry, chemosensors, chemproc, children, chips, cimb, civileng, cleantechnol, climate, clinbioenerg, clinpract, clockssleep, cmd, cmtr, coasts, coatings, colloids, colorants, commodities, complexities, complications, compounds, computation, computers, condensedmatter, conservation, constrmater, cosmetics, covid, crops, cryo, cryptography, crystals, csmf, ctn, culture, curroncol, cyber, dairy, data, ddc, dentistry, dermato, dermatopathology, designs, devices, dhi, diabetology, diagnostics, dietetics, digital, disabilities, diseases, diversity, dna, drones, dynamics, earth, ebj, ecm, ecologies, econometrics, economies, edm, education, eesp, ejihpe, electricity, electrochem, electronicmat, electronics, encyclopedia, endocrines, energies, eng, engproc, entomology, entropy, environments, environremediat, epidemiologia, epigenomes, esa, est, famsci, fermentation, fibers, fintech, fire, fishes, fluids, foods, forecasting, forensicsci, forests, fossstud, foundations, fractalfract, fuels, future, futureinternet, futurepharmacol, futurephys, futuretransp, galaxies, games, gases, gastroent, gastrointestdisord, gastronomy, gels, genealogy, genes, geographies, geohazards, geomatics, geometry, geosciences, geotechnics, geriatrics, germs, glacies, grasses, green, greenhealth, gucdd, hardware, hazardousmatters, healthcare, hearts, hemato, hematolrep, hep, heritage, higheredu, highthroughput, histories, horticulturae, hospitals, humanities, humans, hydrobiology, hydrogen, hydrology, hydropower, hygiene, idr, iic, ijcs, ijem, ijerph, ijfs, ijgi, ijmd, ijms, ijns, ijom, ijpb, ijt, ijtm, ijtpp, ime, immuno, informatics, information, infrastructures, inorganics, insects, instruments, inventions, iot, j, jaestheticmed, jal, jcdd, jcm, jcp, jcrm, jcs, jcto, jdad, jdb, jdream, jemr, jeta, jfb, jfmk, jgbg, jgg, jimaging, jintelligence, jlpea, jmahp, jmmp, jmms, jmp, jmse, jne, jnt, jof, joi, joitmc, joma, jop, joptm, jor, journalmedia, jox, jpbi, jphytomed, jpm, jrfm, jsam, jsan, jtaer, jvd, jzbg, kidneydial, kinasesphosphatases, knowledge, labmed, laboratories, lae, land, languages, laws, life, lights, limnolrev, lipidology, liquids, literature, livers, logics, logistics, lubricants, lymphatics, machines, macromol, magnetism, magnetochemistry, make, marinedrugs, materials, materproc, mathematics, mca, measurements, medicina, medicines, medsci, membranes, merits, metabolites, metals, meteorology, methane, metrics, metrology, micro, microarrays, microbiolres, microelectronics, micromachines, microorganisms, microplastics, microwave, minerals, mining, mmphys, modelling, molbank, molecules, mps, msf, mti, multimedia, muscles, nanoenergyadv, nanomanufacturing, nanomaterials, ncrna, ndt, network, neuroglia, neuroimaging, neurolint, neurosci, nitrogen, notspecified, nursrep, nutraceuticals, nutrients, obesities, occuphealth, oceans, ohbm, onco, optics, oral, organics, organoids, osteology, oxygen, pandemics, parasites, parasitologia, particles, pathogens, pathophysiology, peacestud, pediatrrep, pets, pharmaceuticals, pharmaceutics, pharmacoepidemiology, pharmacy, philosophies, photochem, photonics, phycology, physchem, physics, physiologia, plants, plasma, platforms, pollutants, polymers, polysaccharides, populations, poultry, powders, precipitation, precisoncol, preprints, proceedings, processes, prosthesis, proteomes, psf, psychiatryint, psychoactives, psycholint, publications, purification, quantumrep, quaternary, qubs, radiation, rdt, reactions, realestate, receptors, recycling, regeneration, religions, remotesensing, reports, reprodmed, resources, rheumato, risks, rjpm, robotics, rsee, ruminants, safety, sci, scipharm, sclerosis, seeds, sensors, separations, sexes, shi, signals, sinusitis, siuj, skins, smartcities, sna, societies, socsci, software, soilsystems, solar, solids, spectroscj, sports, standards, stats, std, stratsediment, stresses, surfaces, surgeries, suschem, sustainability, symmetry, synbio, systems, tae, targets, taxonomy, technologies, telecom, test, textiles, thalassrep, therapeutics, thermo, timespace, tomography, tourismhosp, toxics, toxins, tph, transplantology, transportation, traumacare, traumas, tri, tropicalmed, universe, urbansci, uro, vaccines, vehicles, venereology, vetsci, vibration, virtualworlds, viruses, vision, waste, water, welding, wem, wevj, wild, wind, women, world, youth, zoonoticdis

%---------
% article
%---------
% The default type of manuscript is "article", but can be replaced by: 
% abstract, addendum, article, benchmark, book, bookreview, briefcommunication, briefreport, casereport, changes, clinicopathologicalchallenge, comment, commentary, communication, conceptpaper, conferenceproceedings, correction, conferencereport, creative, datadescriptor, discussion, entry, expressionofconcern, extendedabstract, editorial, essay, erratum, fieldguide, hypothesis, interestingimages, letter, meetingreport, monograph, newbookreceived, obituary, opinion, proceedingpaper, projectreport, reply, retraction, review, perspective, protocol, shortnote, studyprotocol, supfile, systematicreview, technicalnote, viewpoint, guidelines, registeredreport, tutorial,  giantsinurology, urologyaroundtheworld
% supfile = supplementary materials

%----------
% submit
%----------
% The class option "submit" will be changed to "accept" by the Editorial Office when the paper is accepted. This will only make changes to the frontpage (e.g., the logo of the journal will get visible), the headings, and the copyright information. Also, line numbering will be removed. Journal info and pagination for accepted papers will also be assigned by the Editorial Office.

%------------------
% moreauthors
%------------------
% If there is only one author the class option oneauthor should be used. Otherwise use the class option moreauthors.

%---------
% pdftex
%---------
% The option pdftex is for use with pdfLaTeX. Remove "pdftex" for (1) compiling with LaTeX & dvi2pdf (if eps figures are used) or for (2) compiling with XeLaTeX.

%=================================================================
% MDPI internal commands - do not modify
\firstpage{1} 
\makeatletter 
\setcounter{page}{\@firstpage} 
\makeatother
\pubvolume{1}
\issuenum{1}
\articlenumber{0}
\pubyear{2026}
\copyrightyear{2026}
%\externaleditor{Firstname Lastname} % More than 1 editor, please add `` and '' before the last editor name
\datereceived{ } 
\daterevised{ } % Comment out if no revised date
\dateaccepted{ } 
\datepublished{ } 
%\datecorrected{} % For corrected papers: "Corrected: XXX" date in the original paper.
%\dateretracted{} % For retracted papers: "Retracted: XXX" date in the original paper.
%\doinum{} % Used for some special journals, like molbank
%\pdfoutput=1 % Uncommented for upload to arXiv.org
%\CorrStatement{yes}  % For updates
%\longauthorlist{yes} % For many authors that exceed the left citation part
%\IsAssociation{yes} % For association journals

%=================================================================
% Add packages and commands here. The following packages are loaded in our class file: fontenc, inputenc, calc, indentfirst, fancyhdr, graphicx, epstopdf, lastpage, ifthen, float, amsmath, amssymb, lineno, setspace, enumitem, mathpazo, booktabs, titlesec, etoolbox, tabto, xcolor, colortbl, soul, multirow, microtype, tikz, totcount, changepage, attrib, upgreek, array, tabularx, pbox, ragged2e, tocloft, marginnote, marginfix, enotez, amsthm, natbib, hyperref, cleveref, scrextend, url, geometry, newfloat, caption, draftwatermark, seqsplit
% cleveref: load \crefname definitions after \begin{document}

% Additional packages not already loaded by the MDPI class
\usepackage[most]{tcolorbox}
% Fix fancyhdr warning about header height
\setlength{\headheight}{20.0pt}

%=================================================================
% Please use the following mathematics environments: Theorem, Lemma, Corollary, Proposition, Characterization, Property, Problem, Example, ExamplesandDefinitions, Hypothesis, Remark, Definition, Notation, Assumption
%% For proofs, please use the proof environment (the amsthm package is loaded by the MDPI class).

%=================================================================
% Full title of the paper (Capitalized)
\Title{UzABSA-LLM: Parameter-Efficient Fine-Tuning and Multi-Domain Annotation for Uzbek Aspect-Based Sentiment Analysis}

% Author Orchid ID: enter ID or remove command
\newcommand{\orcidauthorA}{0000-0002-6895-3436} % Add \orcidA{} behind the author's name
\newcommand{\orcidauthorB}{0000-0002-3664-4954} % Add \orcidB{} behind the author's name
\newcommand{\orcidauthorC}{0000-0002-0369-6707} % 
\newcommand{\orcidauthoD}{0000-0002-7876-5009} % 

% Authors, for the paper (add full first names)
\Author{Sanatbek Matlatipov $^{1}$\orcidA{}, Khabibulla Madatov $^{2}$\orcidB{}, Jaloliddin Rajabov $^{1}$\orcidC{} and Jernej Vi\v{c}i\v{c} $^{3}$\orcidC{}}

%\longauthorlist{yes}

% MDPI internal command: Authors, for metadata in PDF
\AuthorNames{Sanatbek, Khabibulla, Jaloliddin Rajabov and Jernej}

%\longauthorlist{yes}

% Affiliations / Addresses (Add [1] after \address if there is only one affiliation.)
\address{%
$^{1}$ \quad Faculty of Applied Mathematics and Intelligent Technologies, National University of Uzbekistan named after Mirzo Ulugbek, Tashkent 100174, Uzbekistan; {s.matlatipov@nuu.uz}\\
$^{2}$ \quad The Department of Computer Science, Urgench State University Named after Abu Rayhan Biruni,
14 Kh. Alimdjan Str., Urgench City 220100, Uzbekistan; habi1972@mail.ru; \\
$^{3}$ \quad Faculty of Mathematics, Natural Science and Information Technologies, University of Primorska, 6000 Koper, Slovenia; jernej.vicic@upr.si\\
}

% Contact information of the corresponding author
\corres{Correspondence: s.matlatipov@nuu.uz;}

% Current address and/or shared authorship
%\firstnote{Current address: Affiliation.}  
% Current address should not be the same as any items in the Affiliation section.

%\secondnote{These authors contributed equally to this work.}
% The commands \thirdnote{} till \eighthnote{} are available for further notes.

%\simplesumm{} % Simple summary

%\conference{} % An extended version of a conference paper

% Abstract (Do not insert blank lines, i.e. \\) 
\abstract{Aspect-Based Sentiment Analysis (ABSA) remains largely unexplored for Uzbek, a low-resource Turkic language with limited representation in current foundation models. This study presents the first systematic investigation of parameter-efficient fine-tuning for Uzbek ABSA, comparing three open-source large language models: Qwen~2.5-7B, Llama ~3.1-8B, and DeepSeek-7B, each evaluated under identical QLoRA configurations on a SemEval-format restaurant review dataset comprising 6,175 annotated sentences. ABSA is reformulated as a generation task with Uzbek-language instruction tuning, enabling models to jointly extract aspect terms and classify sentiment polarity in a single pass, producing structured JSON output. Evaluation on 609 validation examples demonstrates that Qwen~2.5-7B achieves the highest aspect term extraction F1 score of 0.6603 with a perfect 100\% JSON parse rate, while Llama~3.1-8B leads in end-to-end pair F1 (0.5805) and sentiment accuracy (0.8864). Importantly, training loss does not reliably predict downstream task performance: Qwen surpasses Llama in extraction despite a 31\% higher validation loss. Based on these findings, a three-layer model-assisted annotation pipeline is proposed, integrating local inference, LLM-as-Judge quality scoring, and quality-tiered assembly. Application of this pipeline to 5,038 multi-domain Uzbek business reviews across 23 categories produces a novel silver-standard dataset containing 9,884 aspect–sentiment pairs, with 63.5\% of judged annotations meeting inclusion quality thresholds.}

% Keywords
\keyword{Aspect-Based Sentiment Analysis; Low-Resource NLP; Uzbek Language; Parameter-Efficient Fine-Tuning; QLoRA; Large Language Models; Multi-Domain Dataset} 

% The fields PACS, MSC, and JEL may be left empty or commented out if not applicable
%\PACS{J0101}
%\MSC{}
%\JEL{}

%%%%%%%%%%%%%%%%%%%%%%%%%%%%%%%%%%%%%%%%%%
% Only for the journal Diversity
%\LSID{\url{http://}}

%%%%%%%%%%%%%%%%%%%%%%%%%%%%%%%%%%%%%%%%%%
% Only for the journal Applied Sciences
%\featuredapplication{Authors are encouraged to provide a concise description of the specific application or a potential application of the work. This section is not mandatory.}
%%%%%%%%%%%%%%%%%%%%%%%%%%%%%%%%%%%%%%%%%%

%%%%%%%%%%%%%%%%%%%%%%%%%%%%%%%%%%%%%%%%%%
% Only for the journal Data
%\dataset{DOI number or link to the deposited data set if the data set is published separately. If the data set shall be published as a supplement to this paper, this field will be filled by the journal editors. In this case, please submit the data set as a supplement.}
%\datasetlicense{License under which the data set is made available (CC0, CC-BY, CC-BY-SA, CC-BY-NC, etc.)}

%%%%%%%%%%%%%%%%%%%%%%%%%%%%%%%%%%%%%%%%%%
% Only for the journal BioTech, Fishes, Neuroimaging and Toxins
%\keycontribution{The breakthroughs or highlights of the manuscript. Authors can write one or two sentences to describe the most important part of the paper.}

%%%%%%%%%%%%%%%%%%%%%%%%%%%%%%%%%%%%%%%%%%
% Only for the journal Encyclopedia
%\encyclopediadef{For entry manuscripts only: please provide a brief overview of the entry title instead of an abstract.}

%%%%%%%%%%%%%%%%%%%%%%%%%%%%%%%%%%%%%%%%%%
% Different journals have different requirements. Please check the specific journal guidelines in the "Instructions for Authors" on the journal's official website.
%\addhighlights{yes}
%\renewcommand{\addhighlights}{%
%
%\noindent The goal is to increase the discoverability and readability of the article via search engines and other scholars. Highlights should not be a copy of the abstract, but a simple text allowing the reader to quickly and simplified find out what the article is about and what can be cited from it. Each of these parts should be devoted up to 2~bullet points.\vspace{3pt}\\
%\textbf{What are the main findings?}
% \begin{itemize}[labelsep=2.5mm,topsep=-3pt]
% \item First bullet.
% \item Second bullet.
% \end{itemize}\vspace{3pt}
%\textbf{What are the implications of the main findings?}
% \begin{itemize}[labelsep=2.5mm,topsep=-3pt]
% \item First bullet.
% \item Second bullet.
% \end{itemize}
%}

%%%%%%%%%%%%%%%%%%%%%%%%%%%%%%%%%%%%%%%%%%
\begin{document}
%%%%%%%%%%%%%%%%%%%%%%%%%%%%%%%%%%%%%%%%%%

\section{Introduction}
\label{sec:introduction}
Aspect-Based Sentiment Analysis (ABSA) identifies specific aspects and their polarities in text \cite{pontiki-etal-2014-semeval}. While high-resource languages dominate the field, Uzbek—a Turkic language with almost 38 million speakers\footnote{https://www.worldometers.info/world-population/uzbekistan-population/} which remains less explored due to script duality (Latin/Cyrillic), agglutinative morphology, and limited representation in multilingual pre-training. Recent 7-8B parameter LLMs like Qwen 2.5 \cite{Qwen2024-kr}, Llama 3.1 \cite{Grattafiori2024-ka}, and DeepSeek-R1 \cite{DeepSeek-AI2026-vt} offer cross-lingual potential. Using QLoRA \cite{10.5555/3666122.3666563}, we adapt these models on consumer hardware by reformulating ABSA as a single-pass JSON generation task.

\begin{description}
\item[\textbf{RQ1.}] To what extent can parameter-efficient fine-tuning (QLoRA) of open-source LLMs achieve effective joint aspect extraction and sentiment classification for Uzbek text?

\item[\textbf{RQ2.}] How do different LLM architectures (Qwen~2.5-7B, Llama~3.1-8B, DeepSeek-R1-Distill-7B) compare on Uzbek ABSA under identical QLoRA configurations, and does training loss reliably predict downstream task performance?

\item[\textbf{RQ3.}] Can a model trained exclusively on single-domain data (restaurant reviews) generalize to produce usable ABSA annotations across 23~diverse business domains, and what quality characteristics emerge across domains?
\end{description}

\noindent Our main contributions are as follows:

\begin{enumerate}
\item \textbf{First systematic evaluation} of QLoRA-fine-tuned open-source LLMs for Uzbek ABSA, comparing three architectures at the same parameter scale under fully controlled conditions.

\item \textbf{A scalable model-assisted annotation pipeline} combining local inference, LLM-as-Judge quality scoring via GPT-4o-mini, and quality-tiered dataset assembly---a practical and reproducible methodology for creating ABSA resources in low-resource settings.

\item \textbf{A novel multi-domain silver-standard Uzbek ABSA dataset} covering 5{,}038 reviews across 23~business categories with 9{,}884 annotated aspect--sentiment pairs and per-annotation quality metadata.

\item \textbf{Empirical evidence} that training cross-entropy loss does not reliably predict ABSA task performance: the model with the lowest validation loss (Llama~3.1-8B) is outperformed on aspect extraction by Qwen~2.5-7B, which exhibits 31\% higher loss---a finding with methodological implications for model selection in structured output tasks.
\end{enumerate}
 To support reproducible research in low-resource NLP, all code, datasets, and fine-tuned models are publicly available on GitHub\footnote{\url{https://github.com/SanatbekMatlatipov/UzABSA-LLM}} and HuggingFace\footnote{\url{https://huggingface.co/Sanatbek/UzABSA-LLM}}.


The remainder of this paper is organized as follows. \Cref{sec:related_work} reviews related work on ABSA, parameter-efficient fine-tuning for low-resource languages, and LLM-based evaluation. \Cref{sec:datasets} describes the two datasets used in this study. \Cref{sec:methodology} presents the methodology, including model selection, the QLoRA configuration, and the annotation pipeline. \Cref{sec:results} reports the experimental setup and results. \Cref{sec:discussion} discusses implications and limitations, and \Cref{sec:conclusion} concludes with directions for future work.



%%%%%%%%%%%%%%%%%%%%%%%%%%%%%%%%%%%%%%%%%%

\section{Related Work}
\label{sec:related_work}

% ----------------------------------------------------------------------
SemEval-2014 Task~4~\cite{pontiki-etal-2014-semeval} established the modern ABSA paradigm by defining Aspect Term Extraction (ATE)~\cite{ma-etal-2019-exploring}, Aspect Category Detection (ACD)~\cite{hu-etal-2021-multi-label}, Aspect Sentiment Classification (ASC)~\cite{10.1145/3721844}, and joint end-to-end ABSA on English restaurant and laptop reviews~\cite{Hua2024}. Early systems typically combined sequence labeling (e.g., CRF/BiLSTM) for extraction with attention-based sentiment classifiers. With the rise of pre-trained Transformers, BERT-style models quickly dominated ABSA benchmarks~\cite{ZohrehMadhoushi2019,electronics12030515}, casting ATE as token classification and ASC as sentence(-pair) classification via fine-tuning.

Recent work increasingly reframes ABSA as generation: Scaria et al. \cite{scaria-etal-2024-instructabsa} propose InstructABSA, using instruction-tuned LLMs to jointly extract aspect terms and predict sentiment in a single step by generating structured text, enabling unified modeling and more open-ended aspect vocabularies. However, generative ABSA has largely focused on English, leaving morphologically rich, low-resource languages underexplored; we extend this paradigm to Uzbek by fine-tuning open-source instruction-tuned LLMs to produce structured JSON outputs.% ----------------------------------------------------------------------

\textbf{LLM Fine-Tuning for Low-Resource Languages}. Large multilingual foundation models have broadened NLP beyond English. Architectures such as mT5, BLOOM, Llama~\cite{Touvron2023-lf,Grattafiori2024-ka}, and Qwen~\cite{Qwen2024-kr} differ in low-resource coverage, largely reflecting their pre-training data; for Turkic languages, Qwen~2.5 reports broader multilingual coverage than Llama~3.1, but neither explicitly guarantees Uzbek support. Task adaptation in low-resource settings is constrained by the cost of full fine-tuning for 7--8B models (tens of GB VRAM and substantial compute). Parameter-efficient methods mitigate this: LoRA~\cite{10.5555/3666122.3666563} adds trainable low-rank adapters to attention while freezing base weights, and QLoRA~\cite{10.5555/3666122.3666563} combines LoRA with 4-bit NormalFloat quantization to reduce memory to $<6$~GB for 7B-scale models, enabling single-GPU fine-tuning.
Prior work reports strong results for parameter-efficient adaptation on several non-English/low-resource tasks (e.g., Arabic sentiment analysis, Turkish NER, Korean QA), yet, to our knowledge, LoRA/QLoRA has not been applied to Uzbek NLP. We address this gap by systematically evaluating QLoRA-adapted LLMs for Uzbek ABSA under controlled conditions.% ----------------------------------------------------------------------

\textbf{Uzbek Natural Language Processing}. Early work addressed core infrastructure: Aripov et al. \cite{10.1007/978-3-031-05328-3_6} studied multilingual MT with Uzbek, highlighting challenges from agglutinative morphology, while \cite{sharipov-et-all,sharipov2023uzbektagger} introduced UzbekTagger, a rule-based POS tagger and one of the few public morphological tools. Script variation is another persistent obstacle: after the 1993 shift from Cyrillic to Latin, much user text still appears in Cyrillic; Salaev et al.\cite{Salaev2024} proposed a transliteration tool to normalize across scripts.

For sentiment analysis, Kuriyozov and Matlatipov~\cite{kuriyozov2019building} created an early sentence-level Uzbek dataset with baseline classifiers, later extended by \cite{kuriyozov2022construction} through a broader study of low-resource sentiment dataset construction and benchmarking. Matlatipov et al.\cite{Matlatipov2009,uzbekendingssanat} further examined restaurant-domain sentiment classification on locally collected reviews. The closest resource to our study is UzABSA~\cite{matlatipov2024uzabsa}, the first Uzbek ABSA dataset, following the SemEval-2014 Task~4 scheme~\cite{pontiki-etal-2014-semeval,annoation}, with 6{,}175 annotated restaurant sentences, 7{,}412 aspect-term spans with polarity, and 7{,}724 aspect-category labels over five categories (\textit{ovqat}, \textit{muhit}, \textit{xizmat}, \textit{narx}, \textit{boshqalar}).

Overall, prior Uzbek NLP,including UzABSA\cite{matlatipov2024uzabsa}, has mainly used rule-based, traditional ML, or encoder-based neural models; to our knowledge, no published work has fine-tuned LLMs (full or parameter-efficient) for Uzbek. We address this gap by applying instruction-tuned LLMs with QLoRA-based adaptation to Uzbek NLP.% ----------------------------------------------------------------------

\textbf{LLM-as-Judge for Annotation Quality}. LLM-as-Judge uses large language models as automated evaluators, providing a scalable alternative to human assessment of generated text. Zheng et al.\cite{10.5555/3666122.3668142} introduced MT-Bench and Chatbot Arena, showing that strong LLMs (e.g., GPT-4) can yield scores that correlate with human preferences in open-ended dialogue. This approach has since been applied to instruction-following, summarization, and translation, offering cost efficiency, reproducibility, and scalability when human annotation, especially for low-resource languages which is expensive or impractical. However, LLM judges may suffer from positional bias, self-enhancement bias, and session-to-session score drift\cite{2506.22316}. While these issues are documented for chat-style evaluation, they are less studied for structured tasks such as information extraction or ABSA.

We use LLM-as-Judge for a different goal: assessing structured ABSA annotations produced by a fine-tuned model in a low-resource setting without gold references. Concretely, GPT-4o-mini scores Uzbek aspect--sentiment outputs on five rubric dimensions (completeness, accuracy, sentiment correctness, relevance, overall quality) using a 1--5 Likert scale. Unlike prior work, our target is structured extraction rather than free-form text, and the judge must understand Uzbek input; we then apply stratified sampling over 23 domains and quality thresholds to guide dataset assembly (\Cref{sec:meth_pipeline}).


%%%%%%%%%%%%%%%%%%%%%%%%%%%%%%%%%%%%%%%%%%

%%%%%%%%%%%%%%%%%%%%%%%%%%%%%%%%%%%%%%%%%%

\section{Datasets}
\label{sec:datasets}
This study draws on two complementary corpora. The first is the UzABSA dataset~\cite{matlatipov2024uzabsa}, a publicly available collection of 6{,}175 annotated Uzbek restaurant review sentences hosted on HuggingFace.\footnote{\url{https://huggingface.co/datasets/Sanatbek/aspect-based-sentiment-analysis-uzbek}} After preprocessing and converting the original 6{,}175 examples into ChatML instruction which response instances, we retain 6{,}089 usable training instances, which we split 90/10 into 5{,}480 training and 609 validation examples (seed=42). Each sentence is annotated following the SemEval-2014 Task~4 scheme~\cite{pontiki-etal-2014-semeval}: aspect terms are marked with character-level spans and assigned one of four polarity labels (\textit{positive}, \textit{negative}, \textit{neutral}, \textit{conflict}), and aspect categories are independently labeled with polarity. The dataset contains 7{,}412 aspect term annotations and 7{,}724 aspect category annotations distributed across five domain-specific categories: \textit{ovqat} (food), \textit{muhit} (ambiance), \textit{xizmat} (service), \textit{narx} (price), and \textit{boshqalar} (other). On average, each sentence contains 2.53 aspect annotations and 6.42 words.

A notable characteristic of the dataset is class imbalance: positive polarity accounts for 56\% of term-level annotations (4{,}153 of 7{,}412), followed by neutral (21.6\%), negative (21.0\%), and conflict (1.4\%). This skew toward positive sentiment reflects the natural distribution of restaurant reviews but may bias model predictions toward the majority class. We report per-class metrics alongside micro-averaged scores to quantify this effect (\Cref{sec:results}).  \Cref{tab:dataset_stats} summarizes the key statistics.
\begin{table}[t]
\centering
\caption{Statistics of the annotated UzABSA dataset used for training and evaluation.}
\label{tab:dataset_stats}
\setlength{\tabcolsep}{6pt} % tighten columns (optional)
\begin{tabular}{l r @{\hspace{14pt}} l r}
\toprule
\multicolumn{2}{c}{\textbf{Dataset statistics}} & \multicolumn{2}{c}{\textbf{Term-level polarity distribution}} \\
\midrule
Total examples & 6{,}175 & Positive & 4{,}153 (56.0\%) \\
Aspect term annotations & 7{,}412 & Neutral & 1{,}601 (21.6\%) \\
Aspect category annotations & 7{,}724 & Negative & 1{,}555 (21.0\%) \\
Unique categories & 5 & Conflict & 103 (1.4\%) \\
Avg.\ aspects per sentence & 2.53 & Train / Val split & 5{,}480 / 609 \\
Avg.\ words per sentence & 6.42 & & \\
\bottomrule
\end{tabular}
\end{table}

% ----------------------------------------------------------------------
\subsection{Multi-Domain Uzbek Review Corpus for Model-Assisted Annotation}
\label{sec:data_multidomain}

The second corpus is a novel collection of business reviews scraped from Sharh,\footnote{\url{https://sharh.commeta.uz/en}} a publicly accessible Uzbek-language business review platform. After deduplication and minimum-length filtering (removing entries with fewer than five words), the corpus comprises 5{,}038 reviews spanning 630~unique businesses across 23~business domains. This corpus is not used for training; it serves exclusively as the annotation target for the model-assisted pipeline described in \Cref{sec:meth_pipeline}.

Data collection and usage were conducted with explicit permission from the data owner, as noted in \Cref{sec:acknowledgement}. The annotated corpus derived from these reviews is publicly released on Zenodo to facilitate further research on low-resource ABSA.\footnote{\url{https://doi.org/10.5281/zenodo.18790639}}

The domain distribution is highly skewed: \textit{Restoran/Ovqatlanish} (restaurants and food service) is the largest domain with 1{,}626 reviews (32.3\%), followed by \textit{Bank/Moliya} (banking and finance) with 577 reviews (11.5\%). The ten most frequent domains together account for approximately 89.9\% of the corpus; \Cref{tab:domain_dist} lists their counts and proportions.

\begin{table}[t]
\centering
\caption{Top-10 domain distribution in the multi-domain Uzbek review corpus (5{,}038 reviews total).}
\label{tab:domain_dist}
\begin{tabular}{llrr}
\toprule
\textbf{Rank} & \textbf{Domain} & \textbf{Reviews} & \textbf{\%} \\
\midrule
1  & Restoran/Ovqatlanish (Restaurant/Dining)              & 1{,}626 & 32.3 \\
2  & Bank/Moliya (Banking/Finance)                         &   577   & 11.5 \\
3  & Boshqa (Other)                                        &   507   & 10.1 \\
4  & To'lov tizimlari (Payment systems)                    &   328   &  6.5 \\
5  & Elektron tijorat (E-commerce)                         &   296   &  5.9 \\
6  & Ta'lim (Education)                                    &   287   &  5.7 \\
7  & Tibbiyot/Sog'liqni saqlash (Healthcare/Medicine)      &   285   &  5.7 \\
8  & Telekommunikatsiya (Telecommunications)               &   234   &  4.6 \\
9  & Oziq-ovqat do'konlari (Grocery stores)                &   215   &  4.3 \\
10 & Gul/Sovg'a (Flowers/Gifts)                            &   173   &  3.4 \\
\midrule
   & Remaining 13 domains                                  &   510   & 10.1 \\
   & \textbf{Total}                                        & \textbf{5{,}038} & \textbf{100.0} \\
\bottomrule
\end{tabular}
\end{table}
Text statistics reflect the informal, often brief nature of online business reviews: the average review contains 13.55~words and 96.26~characters, considerably shorter than the restaurant review sentences in the UzABSA training corpus (6.42~words per sentence, which represent clause-level annotations rather than full reviews). A script-level analysis of the corpus reveals that 92.8\% of reviews are written in Uzbek Latin script, 5.2\% in Russian Cyrillic, and 2.0\% in Uzbek Cyrillic---a distribution consistent with the ongoing post-Soviet script transition and the presence of a Russian-speaking minority among Uzbek internet users. No script normalization or transliteration was applied prior to annotation; the fine-tuned model was evaluated on its ability to handle this natural orthographic variation.

% ----------------------------------------------------------------------
\subsection{Task Formulation}
\label{sec:data_task}

We cast ABSA as conditional text generation rather than separate sequence labeling/classification. Given an Uzbek review, the model produces a single structured JSON containing all aspect terms with their categories and sentiment polarities, avoiding multi-stage pipelines, reducing error propagation, and enabling end-to-end optimization.

We use a ChatML-style instruction template with system/user/assistant roles: the system sets Uzbek-ABSA expertise, the user supplies the review and extraction directives, and the assistant returns the JSON annotation. The full formatting is shown in \Cref{fig:instruction_example_ui}.

\begin{figure}[t]
\centering
\begin{minipage}{0.95\textwidth}

\begin{tcolorbox}[enhanced, colback=red!5!white, colframe=red!75!black, 
  title=\texttt{<|im\_start|>system}, fonttitle=\bfseries, 
  attach boxed title to top left={yshift=-2mm, xshift=5mm},
  boxed title style={colback=red!75!black}]
\footnotesize
Siz o'zbek tilida matnlardan aspektlarni va ularning hissiyotlarini aniqlash bo'yicha mutaxassissiz.\\
Berilgan matndan barcha aspekt terminlarini, ularning kategoriyalarini va hissiyot polaritesini (positive, negative, neutral) ajratib oling.\\
Javobni quyidagi Python dictionary formatida qaytaring:\\
\texttt{\{ "aspects": [ \{ "term": "...", "category": "...", "polarity": "..." \} ] \}}
\end{tcolorbox}

\vspace{-2mm}

\begin{tcolorbox}[enhanced, colback=blue!5!white, colframe=blue!75!black, 
  title=\texttt{<|im\_start|>user}, fonttitle=\bfseries, 
  attach boxed title to top left={yshift=-2mm, xshift=5mm},
  boxed title style={colback=blue!75!black}]
\footnotesize
Quyidagi o'zbek tilidagi matndan aspektlarni, kategoriyalarni va sentimentini aniqlang:\\
\textbf{Matn:} "Bu restoranning ovqatlari juda mazali, lekin narxlari qimmat."
\end{tcolorbox}

\vspace{-2mm}

\begin{tcolorbox}[enhanced, colback=green!5!white, colframe=green!50!black, 
  title=\texttt{<|im\_start|>assistant}, fonttitle=\bfseries, 
  attach boxed title to top left={yshift=-2mm, xshift=5mm},
  boxed title style={colback=green!50!black}]
\footnotesize\ttfamily
\{"aspects": [\\
\ \ \{"term": "ovqatlari", "category": "ovqat", "polarity": "positive"\},\\
\ \ \{"term": "narxlari", "category": "narx", "polarity": "negative"\}\\
]\}
\end{tcolorbox}

\end{minipage}
\caption{Complete ChatML instruction template for Uzbek ABSA. The example demonstrates joint extraction of aspect terms (\textit{ovqatlari} ``food'', \textit{narxlari} ``prices''), category assignment (\textit{ovqat}, \textit{narx}), and sentiment classification (positive, negative) in a single generation step.}
\label{fig:instruction_example_ui}
\end{figure}

The JSON output format serves multiple methodological purposes: (a)~it ensures structured, parseable responses that enable automatic metric computation without post-processing; (b)~it naturally accommodates variable numbers of aspects per review; and (c)~it facilitates batch annotation pipelines for large-scale corpus creation. Each aspect annotation contains three fields: \texttt{term} (the extracted span), \texttt{category} (one of the five UzABSA categories: \textit{ovqat}, \textit{muhit}, \textit{xizmat}, \textit{narx}, \textit{boshqalar}), and \texttt{polarity} (positive, negative, neutral, or conflict).

This generative formulation allows the model to leverage its pre-trained language understanding to simultaneously identify aspect boundaries, assign semantic categories, and classify sentiment polarity---tasks that traditional approaches handle sequentially. The instruction text is provided entirely in Uzbek to maintain language consistency and reduce code-switching effects that could degrade model performance on this low-resource language.



%%%%%%%%%%%%%%%%%%%%%%%%%%%%%%%%%%%%%%%%%%

\section{Methodology}
\label{sec:methodology}

% ----------------------------------------------------------------------
Our design requires models with (i) open-source, permissively licensed weights, (ii) instruction-tuned variants for complex prompt following, (iii) $\sim$7--8B parameters for controlled comparison under a fixed compute budget, and (iv) 4-bit quantized availability compatible with QLoRA. Under these constraints, we select three complementary LLMs: Qwen~2.5-7B-Instruct~\cite{Qwen2024-kr}, Llama~3.1-8B-Instruct~\cite{Grattafiori2024-ka}, and DeepSeek-R1-Distill-Qwen-7B (a Qwen-7B backbone distilled from DeepSeek-R1~\cite{DeepSeek-AI2026-vt}).

Qwen~2.5-7B-Instruct~\cite{Qwen2024-kr} is a 7.6B decoder-only transformer with multilingual pre-training (reported $>$29 languages) and a GQA+SwiGLU+RoPE configuration (28 heads, 4 KV heads), which may benefit Uzbek. Llama~3.1-8B-Instruct~\cite{Grattafiori2024-ka} is an 8B GQA+SwiGLU+RoPE model (32 heads, 8 KV heads) with different depth/width (32 layers; 4{,}096 hidden vs.\ Qwen's 28 layers; 3{,}584), and its Turkic coverage is less explicitly documented. DeepSeek-R1-Distill-Qwen-7B matches Qwen architecturally but differs in weights due to reasoning-oriented distillation rather than standard instruction tuning~\cite{DeepSeek-AI2026-vt}. This choice supports two controlled contrasts: Qwen vs.\ Llama isolates architectural effects under instruction tuning, while Qwen vs.\ Distill-Qwen isolates the pre-training objective on a fixed architecture. All models are loaded from HuggingFace Hub in pre-quantized 4-bit NormalFloat via Unsloth,\footnote{\url{https://github.com/unslothai/unsloth}} avoiding quantization as a confound.
% ----------------------------------------------------------------------
\subsection{QLoRA Fine-Tuning}
\label{sec:meth_qlora}

We adopt QLoRA~\cite{10.5555/3666122.3666563} for parameter-efficient adaptation. Base weights are quantized to 4-bit NormalFloat (NF4) with double quantization and BFloat16 compute. LoRA adapters ($r{=}16$, $\alpha{=}32$, dropout $0.05$, no bias) are applied to all attention and feed-forward linear projections, yielding 40.4M trainable parameters ($\sim$0.92\% of total for the 7B models). Training is performed via the SFTTrainer with Unsloth's optimized fused kernels and gradient checkpointing to minimize memory footprint. All three models are trained under strictly identical hyperparameters to isolate architectural differences: 8-bit AdamW optimizer, learning rate of $2{\times}10^{-4}$ with a cosine schedule and 0.1 warmup ratio, 0.01 weight decay, and gradient clipping at 1.0. We train for 1{,}000 steps on the 5{,}480-example split using a maximum sequence length of 2{,}048 tokens and an effective batch size of 16 (per-device batch 4, 4 accumulation steps).
% ----------------------------------------------------------------------
\subsection{Multi-Domain Annotation Pipeline}
\label{sec:meth_pipeline}

To build a silver-standard, multi-domain Uzbek ABSA dataset, we design a three-layer pipeline that combines local inference with external quality validation. This addresses a common low-resource constraint: for the 23 out-of-domain business categories, no gold annotations exist and large-scale Uzbek expert labeling is prohibitively costly.

\textbf{Layer~1: Model Inference.} We select the best-performing model from the three-model comparison (\Cref{sec:res_comparison}) as the annotation engine, prioritising operational reliability: 100\% JSON parse rate, highest ATE recall (maximising aspect coverage), and fastest inference throughput. Under these criteria, Qwen~2.5-7B is chosen (\Cref{sec:res_comparison}). We run greedy batch inference over all 5{,}038 reviews using the same instruction template and hyperparameters as evaluation, with checkpoint-based crash recovery (state saved every 100 reviews).

\textbf{Layer~2: LLM-as-Judge Quality Scoring.} Without references, we use GPT-4o-mini (via the OpenAI API) as an external judge. We draw a stratified sample of 307 reviews across all 23 domains to cover both high-frequency domains (e.g., restaurants, $n{=}50$) and low-frequency domains (e.g., insurance, $n{=}3$). For each sampled review, the judge assigns 1--5 Likert scores for five rubric dimensions: completeness (coverage of expressed opinions), accuracy (terms present in the text), sentiment correctness (polarity labels), relevance (category--domain fit), and overall quality. The judge relies only on Uzbek text comprehension, consistent with the reference-free LLM-as-Judge paradigm in \Cref{sec:rw_judge}.

\textbf{Layer~3: Quality-Tiered Assembly.} We partition outputs into \emph{Include} (overall $\geq 3.5$), \emph{Flag for review} ($2.5 \leq$ overall $< 3.5$), and \emph{Exclude} (overall $< 2.5$). The assembly script joins predictions with quality metadata (per-dimension scores, business category, source URL) and exports: (i) full JSON with provenance metadata, (ii) a filtered silver-standard subset (included + unjudged reviews), and (iii) JSONL for direct loading via HuggingFace \texttt{datasets}. Non-sampled reviews are released with an ``unjudged'' flag; the 63.5\% inclusion rate on the judged subset provides a statistical basis for their expected quality. The resulting corpus is released on Zenodo (DOI: \href{https://doi.org/10.5281/zenodo.18790639}{10.5281/zenodo.18790639}).
% ----------------------------------------------------------------------
\subsection{Evaluation Framework}
\label{sec:meth_eval}

We evaluate model performance along four complementary dimensions, each targeting a distinct facet of the ABSA generation task.

\textbf{Aspect Term Extraction (ATE).} We report micro-averaged precision, recall, and F1 under two matching criteria. \emph{Exact match} requires the predicted term to be identical to the gold term after case-folding and whitespace normalisation. \emph{Partial match} relaxes this constraint: a prediction is counted as correct if the predicted string is a substring of the gold term or vice versa (i.e., a superstring relationship), accommodating morphological variation in Uzbek---for example, matching \textit{ovqatlari} against \textit{ovqat}. In both cases, metrics are aggregated across all test examples (micro-averaging) to avoid bias from examples with few or many aspect terms.

\textbf{End-to-End ABSA (E2E).} The primary evaluation metric is the \emph{aspect--polarity pair F1}, which requires both the aspect term and its sentiment polarity to be correct for a true positive. This metric captures joint extraction and classification performance in a single score and is directly comparable across models.

\textbf{Sentiment Classification.} For aspect terms that are correctly matched (exact match), we additionally report sentiment accuracy and macro-averaged F1 over the polarity classes. This isolates the model's sentiment assignment capability from its term extraction performance.

\textbf{Instruction Following.} We measure the \emph{JSON parse success rate}---the percentage of model outputs that yield a syntactically valid JSON object containing the expected \texttt{aspects} key. Failed parses are excluded from ABSA metric computation; a low parse rate therefore both reduces the effective evaluation set and signals poor instruction compliance.

All inference is performed with near-greedy decoding (\texttt{temperature}$=$0.1, \texttt{do\_sample}$=$\texttt{False}, \texttt{max\_new\_tokens}$=$512) to ensure deterministic, reproducible outputs. For the multi-domain annotation quality assessment, we supplement automatic metrics with the five-dimension LLM-as-Judge rubric described in \Cref{sec:meth_pipeline}.


%%%%%%%%%%%%%%%%%%%%%%%%%%%%%%%%%%%%%%%%%%

\section{Results}
\label{sec:results}

% ----------------------------------------------------------------------
All experiments run on a single NVIDIA RTX A6000 GPU (48\,GB VRAM) using Python~3.10 with Unsloth, TRL, PEFT, and bitsandbytes. The UzABSA dataset is split 90/10 into training (5{,}480) and validation (609) sets with a fixed seed of~42. Each of the three models is trained for exactly 1{,}000 gradient steps under the QLoRA configuration described in \Cref{sec:meth_qlora}; the \emph{only} variable across runs is the base model. No zero-shot or encoder-based baselines are included; we reserve such comparisons for future work. \Cref{tab:train_efficiency} summarises the resulting training efficiency.

\begin{table}[t]
\centering
\caption{Training efficiency on a single RTX A6000. All runs use identical QLoRA settings (40.4\,M trainable parameters, 0.92\% of total).}
\label{tab:train_efficiency}
\setlength{\tabcolsep}{4pt}
\begin{tabular}{lccc}
\toprule
\textbf{Metric} & \textbf{Qwen 2.5-7B} & \textbf{Llama 3.1-8B} & \textbf{DeepSeek-R1-7B} \\
\midrule
Compute time (min)       & 49   & 60   & 62   \\
Throughput (samples/s)   & 5.44 & 4.46 & 4.28 \\
GPU mem.\ allocated (GB) & 5.4  & 5.6  & 5.5  \\
Total FLOPs              & $1.24\!\times\!10^{17}$ & $1.50\!\times\!10^{17}$ & $1.42\!\times\!10^{17}$ \\
Best eval loss           & 0.3656 & 0.2785 & 0.3363 \\
\bottomrule
\end{tabular}
\end{table}

% ----------------------------------------------------------------------
\subsection{Three-Model ABSA Comparison (RQ1, RQ2)}
\label{sec:res_comparison}

\Cref{tab:absa_results} presents the full evaluation on the 609-example validation set. Three complementary findings emerge.

\begin{table}[t]
\centering
\caption{ABSA evaluation on the validation set (609 examples). Best scores per metric are \textbf{bold}. All models use greedy decoding (\texttt{temperature}$=$0.1).}
\label{tab:absa_results}
\setlength{\tabcolsep}{3.5pt}
\begin{tabular}{llccc}
\toprule
\textbf{Task} & \textbf{Metric} & \textbf{Qwen 2.5-7B} & \textbf{Llama 3.1-8B} & \textbf{DeepSeek-R1-7B} \\
\midrule
\multirow{3}{*}{ATE (exact)}
  & Precision & 0.6138 & \textbf{0.6394} & 0.5348 \\
  & Recall    & \textbf{0.7143} & 0.6713 & 0.6921 \\
  & F1        & \textbf{0.6603} & 0.6549 & 0.6034 \\
\midrule
\multirow{3}{*}{ATE (partial)}
  & Precision & 0.7163 & \textbf{0.7411} & 0.6452 \\
  & Recall    & \textbf{0.8336} & 0.7781 & 0.8350 \\
  & F1        & \textbf{0.7705} & 0.7591 & 0.7279 \\
\midrule
\multirow{3}{*}{E2E Pair}
  & Precision & 0.5387 & \textbf{0.5667} & 0.4448 \\
  & Recall    & \textbf{0.6269} & 0.5950 & 0.5756 \\
  & F1        & 0.5795 & \textbf{0.5805} & 0.5018 \\
\midrule
\multirow{2}{*}{Sentiment}
  & Accuracy  & 0.8777 & \textbf{0.8864} & 0.8317 \\
  & Macro-F1  & 0.8113 & \textbf{0.8435} & 0.7717 \\
\midrule
JSON parse & Success rate & \textbf{100.0\%} & 95.9\% & 95.4\% \\
\bottomrule
\end{tabular}
\end{table}

\textbf{Qwen excels at aspect extraction.} Qwen~2.5-7B achieves the highest ATE F1 on both exact (0.6603) and partial match (0.7705), driven by the best recall across all extraction metrics. It also attains a perfect 100\% JSON parse rate---a critical advantage for production annotation pipelines.

\textbf{Llama leads on sentiment and end-to-end performance.} Llama~3.1-8B obtains the highest E2E Pair F1 (0.5805), sentiment accuracy (0.8864), and sentiment macro-F1 (0.8435). Its precision-oriented extraction compensates for slightly lower recall, yielding the best aspect--polarity pair matching. Notably, the Pair F1 gap between Llama and Qwen is only 0.001---effectively a tie.

\textbf{DeepSeek underperforms despite reasoning distillation.} DeepSeek-R1-Distill-Qwen-7B scores lowest on every ABSA metric. It over-predicts aspects substantially (933 predictions vs.\ 721 gold; ratio 1.29), whereas Llama is closest to the gold count (757; ratio 1.05) and Qwen moderately over-predicts (839; ratio 1.16). The R1 reasoning distillation, while designed for chain-of-thought tasks, appears detrimental to direct structured JSON extraction.

\textbf{Training Loss vs.\ Task Performance.} A notable finding for RQ2 is that ranking by validation loss, as visualized in Fig. \ref{fig:loss_curves}, diverges from ranking by downstream task metrics.
Llama achieves the lowest evaluation loss (0.2785) and the highest E2E Pair F1 (0.5805), yet Qwen surpasses it on exact ATE F1 (0.6603 vs.\ 0.6549) despite having a 31\% higher validation loss (0.3656). Furthermore, DeepSeek holds the second-best evaluation loss (0.3363) but yields the worst ABSA scores across the board (ATE F1 0.6034; Pair F1 0.5018). Cross-entropy loss captures token-level prediction quality but is blind to structured output correctness---aspect boundary precision, JSON formatting, and polarity alignment. This finding implies that model selection for structured extraction tasks should not rely on validation loss alone.

\begin{figure}[H]
\begin{adjustwidth}{-\extralength}{0cm}
\centering
\includegraphics[width=1.0\linewidth]{wand-uzabsa.png}
\end{adjustwidth}
\caption{Training dynamics of the three fine-tuned LLMs over 1,000 global steps. 
(\textbf{a}) Validation loss curves for Qwen~2.5-7B, Llama~3.1-8B, and DeepSeek-R1-Distill-Qwen-7B. 
(\textbf{b}) Learning-rate schedule used during fine-tuning. 
(\textbf{c}) Training loss curves for the three models.}
\label{fig:loss_curves}
\end{figure}
% ----------------------------------------------------------------------
\subsection{Multi-Domain Quality Assessment (RQ3)}
\label{sec:res_multidomain}

To evaluate cross-domain generalisation, we apply the annotation pipeline (\Cref{sec:meth_pipeline}) using the Qwen model---selected for its 100\% parse rate and highest ATE recall---and score a stratified sample of 307~reviews via GPT-4o-mini as LLM-as-Judge. \Cref{tab:judge_agg} reports the aggregate quality scores across five dimensions.

\begin{table}[t]
\centering
\caption{LLM-as-Judge aggregate quality scores (307-review stratified sample, 1--5 scale).}
\label{tab:judge_agg}
\begin{tabular}{lc l}
\toprule
\textbf{Dimension} & \textbf{Mean} & \textbf{Interpretation} \\
\midrule
Accuracy     & 4.47 & Extracted terms are genuinely present \\
Sentiment    & 4.21 & Polarity labels largely correct \\
Relevance    & 4.06 & Categories contextually appropriate \\
Overall      & 3.75 & Majority suitable for inclusion \\
Completeness & 3.32 & Main opinions captured; some missed \\
\bottomrule
\end{tabular}
\end{table}

Quality-tier assignment using the overall score yields 63.5\% \emph{Include} ($\geq$3.5; $n$=195), 28.0\% \emph{Flag} (2.5--3.49; $n$=86), and 8.5\% \emph{Exclude} ($<$2.5; $n$=26). The high accuracy score (4.47) confirms that the model rarely hallucinates aspects; the lower completeness score (3.32) reflects a tendency to miss subsidiary or domain-specific opinions.

\textbf{Domain-Level Quality Gradient.} \Cref{tab:domain_quality} presents per-domain overall scores for selected domains spanning the quality spectrum. A clear gradient emerges from in-domain to out-of-domain categories: the model trained exclusively on restaurant reviews achieves its highest quality on \emph{Restoran} (4.24) and near-domain categories sharing overlapping aspect vocabularies (e.g., \emph{Sayohat/Turizm}, 4.40), while performance degrades on domains with specialised terminology such as \emph{Telekommunikatsiya} (3.32) and \emph{Ta'lim} (3.25). The most distant domains---\emph{Sug'urta} (insurance) and \emph{Investitsiya} (investment)---drop to~2.67, driven by completeness failures rather than accuracy: the accuracy sub-score remains $\geq$3.67 across all 23~domains. Overall, 15 of 23 domains (65\%) meet the inclusion threshold ($\geq$3.5).

\begin{table}[t]
\centering
\caption{Per-domain LLM-as-Judge quality scores (overall, 1--5 scale). Domains grouped by proximity to the training distribution.}
\label{tab:domain_quality}
\setlength{\tabcolsep}{4pt}
\begin{tabular}{llcccccc}
\toprule
\textbf{Group} & \textbf{Domain} & $N$ & \textbf{Comp.} & \textbf{Acc.} & \textbf{Sent.} & \textbf{Rel.} & \textbf{Overall} \\
\midrule
In-domain    & Restoran       & 50 & 3.68 & 4.76 & 4.56 & 4.66 & 4.24 \\
\midrule
\multirow{2}{*}{Near}
             & Sayohat/Turizm & 5  & 4.00 & 4.80 & 4.60 & 4.40 & 4.40 \\
             & Oziq-ovqat     & 15 & 3.47 & 4.33 & 4.47 & 4.13 & 3.80 \\
\midrule
\multirow{3}{*}{Out-of-domain}
             & Bank/Moliya    & 30 & 3.40 & 4.47 & 4.17 & 3.90 & 3.70 \\
             & Telekom.       & 25 & 3.04 & 4.36 & 3.64 & 3.64 & 3.32 \\
             & Ta'lim         & 20 & 2.90 & 4.20 & 3.80 & 3.45 & 3.25 \\
\midrule
\multirow{2}{*}{Distant}
             & Sug'urta       & 3  & 2.33 & 3.67 & 2.67 & 3.33 & 2.67 \\
             & Investitsiya   & 3  & 2.33 & 4.00 & 3.67 & 3.00 & 2.67 \\
\bottomrule
\end{tabular}
\end{table}

\textbf{Annotation Pipeline Output.} Batch annotation of the full 5{,}038-review corpus produces 9{,}884 aspects (average 1.96 per review) with a 100\% JSON parse rate across all 23~domains. The polarity distribution is 75.1\% positive, 16.5\% negative, and 8.4\% neutral---reflecting a natural skew in user-generated reviews. The aspect category distribution reveals a notable pattern: \emph{boshqalar} (miscellaneous) dominates at 45.5\%, followed by \emph{ovqat} (17.9\%), \emph{xizmat} (15.9\%), \emph{muhit} (10.7\%), and \emph{narx} (9.0\%).


%%%%%%%%%%%%%%%%%%%%%%%%%%%%%%%%%%%%%%%%%%

\section{Discussion}
\label{sec:discussion}

% ----------------------------------------------------------------------
\textbf{The Completeness Bottleneck}. Completeness (mean 3.32/5) is the weakest quality dimension: the model reliably extracts aspects present in the text but often misses domain-specific or subtle opinions, especially outside the restaurant domain. This is primarily due to the training data's narrow coverage—only five aspect categories—forcing the model to map unfamiliar aspects to the generic "boshqalar" (miscellaneous) label, which dominates at 45.5\%. Addressing this bottleneck will require domain-adaptive fine-tuning or expanding the aspect taxonomy to better capture out-of-domain phenomena.

% ----------------------------------------------------------------------
\textbf{Architecture vs.\ Pre-Training Effects}. DeepSeek's R1 reasoning distillation does not benefit structured ABSA: despite sharing the Qwen-7B backbone, it underperforms both Qwen and Llama, likely due to reasoning-oriented output patterns conflicting with strict JSON formatting. Llama achieves the lowest loss and best sentiment accuracy, but Qwen is superior for reliable extraction and instruction following (100\% parse rate, highest ATE recall). For annotation pipelines, robust JSON compliance and extraction coverage are more critical than marginal gains in sentiment metrics.

% ----------------------------------------------------------------------
\textbf{Loss as a Model Selection Proxy}. Validation loss ranking diverges from actual ABSA performance: Llama's lowest loss does not translate to best extraction, while Qwen outperforms on ATE F1 despite higher loss. Cross-entropy loss fails to capture structured output quality, such as JSON validity or aspect boundary precision. Model selection for ABSA and similar tasks should always rely on task-specific metrics, not loss alone.

% ----------------------------------------------------------------------
\textbf{Limitations}. This study is limited by training exclusively on restaurant reviews, with no domain-adaptive fine-tuning or human gold-standard for the multi-domain corpus. No encoder-based baselines (e.g., BERT) are included, and positive polarity dominates both datasets. Tokenizer efficiency and further human verification remain for future work.

%%%%%%%%%%%%%%%%%%%%%%%%%%%%%%%%%%%%%%%%%%


%%%%%%%%%%%%%%%%%%%%%%%%%%%%%%%%%%%%%%%%%%
\section{Conclusions}
\label{sec:conclusion}
This paper presented the first systematic study of parameter-efficient fine-tuning of open-source LLMs for Uzbek ABSA. Addressing \textbf{RQ1}, we demonstrated that QLoRA successfully enables effective joint aspect extraction and sentiment classification via structured JSON generation. For \textbf{RQ2}, our comparison revealed a clear architectural trade-off: Qwen~2.5-7B excels at reliable aspect extraction and instruction following, whereas Llama~3.1-8B achieves the highest precision for end-to-end sentiment matching. Crucially, we found that validation loss is an unreliable proxy for downstream structured extraction performance. Finally, addressing \textbf{RQ3}, we showed that a model trained exclusively on single-domain data can generalize to out-of-domain text. Our proposed three-layer pipeline successfully produced a high-fidelity, multi-domain silver-standard resource, paving the way for scalable low-resource dataset creation. Future work will focus on expanding the aspect taxonomy and exploring domain-adaptive fine-tuning to improve extraction completeness across diverse business domains.
%%%%%%%%%%%%%%%%%%%%%%%%%%%%%%%%%%%%%%%%%%
\vspace{6pt} 

%%%%%%%%%%%%%%%%%%%%%%%%%%%%%%%%%%%%%%%%%%
%% optional
%\supplementary{The following supporting information can be downloaded at:  \linksupplementary{s1}, Figure S1: title; Table S1: title; Video S1: title.}

% Only for journal Methods and Protocols:
% If you wish to submit a video article, please do so with any other supplementary material.
% \supplementary{The following supporting information can be downloaded at: \linksupplementary{s1}, Figure S1: title; Table S1: title; Video S1: title. A supporting video article is available at doi: link.}

% Only used for preprtints:
% \supplementary{The following supporting information can be downloaded at the website of this paper posted on \href{https://www.preprints.org/}{Preprints.org}.}

% Only for journal Hardware:
% If you wish to submit a video article, please do so with any other supplementary material.
% \supplementary{The following supporting information can be downloaded at: \linksupplementary{s1}, Figure S1: title; Table S1: title; Video S1: title.\vspace{6pt}\\
%\begin{tabularx}{\textwidth}{lll}
%\toprule
%\textbf{Name} & \textbf{Type} & \textbf{Description} \\
%\midrule
%S1 & Python script (.py) & Script of python source code used in XX \\
%S2 & Text (.txt) & Script of modelling code used to make Figure X \\
%S3 & Text (.txt) & Raw data from experiment X \\
%S4 & Video (.mp4) & Video demonstrating the hardware in use \\
%... & ... & ... \\
%\bottomrule
%\end{tabularx}
%}

%%%%%%%%%%%%%%%%%%%%%%%%%%%%%%%%%%%%%%%%%%
\authorcontributions{Conceptualization, S.M. and J.V.; methodology, S.M and K.M.; software, S.M. and J.R.; validation, J.V.; investigation, S.M., K.M. and J.R.; resources, S.M.; data curation, J.R.; writing—original
draft preparation, S.M. and K.M.; writing—review and editing, K.M. and J.V.; visualization, S.M.;
supervision, J.V.; project administration, K.M.; funding acquisition, not applicable. All authors have
read and agreed to the published version of the manuscript.}

\funding{This research received no external funding}

\institutionalreview{Not applicable}

\dataavailability{The primary Uzbek aspect-based sentiment analysis dataset used for model fine-tuning is publicly available on Hugging Face at \url{https://huggingface.co/datasets/Sanatbek/aspect-based-sentiment-analysis-uzbek}. The multi-domain Uzbek business review corpus generated in this study is publicly available on Zenodo at \url{https://zenodo.org/records/18790639} with DOI: \href{https://doi.org/10.5281/zenodo.18790639}{10.5281/zenodo.18790639}. The fine-tuned model checkpoints and related resources are available on Hugging Face at \url{https://huggingface.co/Sanatbek/UzABSA-LLM}, and the implementation code is available at \url{https://github.com/SanatbekMatlatipov/UzABSA-LLM}.}

%\durcstatement{Current research is limited to the [please insert a specific academic field, e.g., XXX], which is beneficial [share benefits and/or primary use] and does not pose a threat to public health or national security. Authors acknowledge the dual-use potential of the research involving xxx and confirm that all necessary precautions have been taken to prevent potential misuse. As an ethical responsibility, authors strictly adhere to relevant national and international laws about DURC. Authors advocate for responsible deployment, ethical considerations, regulatory compliance, and transparent reporting to mitigate misuse risks and foster beneficial outcomes.}

% Only for journal Nursing Reports
%\publicinvolvement{Please describe how the public (patients, consumers, carers) were involved in the research. Consider reporting against the GRIPP2 (Guidance for Reporting Involvement of Patients and the Public) checklist. If the public were not involved in any aspect of the research add: ``No public involvement in any aspect of this research''.}
%
%% Only for journal Nursing Reports
%\guidelinesstandards{Please add a statement indicating which reporting guideline was used when drafting the report. For example, ``This manuscript was drafted against the XXX (the full name of reporting guidelines and citation) for XXX (type of research) research''. A complete list of reporting guidelines can be accessed via the equator network: \url{https://www.equator-network.org/}.}
%
%% Only for journal Nursing Reports
%\useofartificialintelligence{Please describe in detail any and all uses of artificial intelligence (AI) or AI-assisted tools used in the preparation of the manuscript. This may include, but is not limited to, language translation, language editing and grammar, or generating text. Alternatively, please state that “AI or AI-assisted tools were not used in drafting any aspect of this manuscript”.}

\acknowledgments{\label{sec:acknowledgement}We extend our heartfelt gratitude to Sardor Berdiyev, CEO and Co-founder of Commeta, for graciously granting permission to use publicly available reviews from the sharh.commeta.uz platform exclusively for academic research purposes. We also acknowledge the computational resources provided by the National University of Uzbekistan. “During the preparation of this manuscript/study, the author Sanatbek used OpenAI's GPT model for the purposes of paraphrasing and better writing purposes. The authors have reviewed and edited the output and take full responsibility for the content of this publication.”}


\conflictsofinterest{ The authors declare no conflicts of interest.} 

%%%%%%%%%%%%%%%%%%%%%%%%%%%%%%%%%%%%%%%%%%
%% Optional

%% Only for journal Encyclopedia
%\entrylink{The Link to this entry published on the encyclopedia platform.}

\abbreviations{Abbreviations}{
The following abbreviations are used in this manuscript:
\\

\noindent 
\begin{tabular}{@{}ll}
ABSA & Aspect-Based Sentiment Analysis\\
ACD & Aspect Category Detection\\
ASC & Aspect Sentiment Classification\\
ATE & Aspect Term Extraction\\
NLP & Natural Language Processing\\
QLoRA & Quantized Low-Rank Adaptation\\
RQ & Research Question
\end{tabular}
}

%%%%%%%%%%%%%%%%%%%%%%%%%%%%%%%%%%%%%%%%%%
%% Optional
\appendixtitles{no} % Leave argument "no" if all appendix headings stay EMPTY (then no dot is printed after "Appendix A"). If the appendix sections contain a heading then change the argument to "yes".
\appendixstart
\appendix
\section[\appendixname~\thesection]{}
\subsection[\appendixname~\thesubsection]{}
The appendix is an optional section that can contain details and data supplemental to the main text---for example, explanations of experimental details that would disrupt the flow of the main text but nonetheless remain crucial to understanding and reproducing the research shown; figures of replicates for experiments of which representative data are shown in the main text can be added here if brief, or as Supplementary Data. Mathematical proofs of results not central to the paper can be added as an appendix.

\begin{table}[H] 
\caption{This is a table caption.\label{tab5}}
%\newcolumntype{C}{>{\centering\arraybackslash}X}
\begin{tabularx}{\textwidth}{CCC}
\toprule
\textbf{Title 1}	& \textbf{Title 2}	& \textbf{Title 3}\\
\midrule
Entry 1		& Data			& Data\\
Entry 2		& Data			& Data\\
\bottomrule
\end{tabularx}
\end{table}

\section[\appendixname~\thesection]{}
All appendix sections must be cited in the main text. In the appendices, Figures, Tables, etc. should be labeled, starting with ``A''---e.g., Figure A1, Figure A2, etc.

%%%%%%%%%%%%%%%%%%%%%%%%%%%%%%%%%%%%%%%%%%
%\isPreprints{}{% This command is only used for ``preprints''.
\begin{adjustwidth}{-\extralength}{0cm}
%} % If the paper is ``preprints'', please uncomment this parenthesis.
%\printendnotes[custom] % Un-comment to print a list of endnotes

\reftitle{References}

% Please provide either the correct journal abbreviation (e.g. according to the “List of Title Word Abbreviations” http://www.issn.org/services/online-services/access-to-the-ltwa/) or the full name of the journal.
% Citations and References in Supplementary files are permitted provided that they also appear in the reference list here. 

%=====================================
% References, variant A: external bibliography
%=====================================
\bibliography{mybibliography}

%=====================================
% References, variant B: internal bibliography
%=====================================

% ACS format
% \isAPAandChicago{}{%
% \begin{thebibliography}{999}
% % Reference 1
% \bibitem[Author1(year)]{ref-journal}
% Author~1, T. The title of the cited article. {\em Journal Abbreviation} {\bf 2008}, {\em 10}, 142--149.
% % Reference 2
% \bibitem[Author2(year)]{ref-book1}
% Author~2, L. The title of the cited contribution. In {\em The Book Title}; Editor 1, F., Editor 2, A., Eds.; Publishing House: City, Country, 2007; pp. 32--58.
% % Reference 3
% \bibitem[Author1 and Author2 (year)]{ref-book2}
% Author 1, A.; Author 2, B. \textit{Book Title}, 3rd ed.; Publisher: Publisher Location, Country, 2008; pp. 154--196.
% % Reference 4
% \bibitem[Author4(year)]{ref-unpublish}
% Author 1, A.B.; Author 2, C. Title of Unpublished Work. \textit{Abbreviated Journal Name} year, \textit{phrase indicating stage of publication (submitted; accepted; in press)}.
% % Reference 5
% \bibitem[Author8(year)]{ref-url}
% Title of Site. Available online: URL (accessed on Day Month Year).
% % Reference 6
% \bibitem[Author6(year)]{ref-proceeding}
% Author 1, A.B.; Author 2, C.D.; Author 3, E.F. Title of presentation. In Proceedings of the Name of the Conference, Location of Conference, Country, Date of Conference (Day Month Year); Abstract Number (optional), Pagination (optional).
% % Reference 7
% \bibitem[Author7(year)]{ref-thesis}
% Author 1, A.B. Title of Thesis. Level of Thesis, Degree-Granting University, Location of University, Date of Completion.
% \end{thebibliography}
% }

% % Chicago format (Used for journal: arts, genealogy, histories, humanities, jintelligence, laws, literature, religions, risks, socsci)
% \isChicagoStyle{%
% \begin{thebibliography}{999}
% % Reference 1
% \bibitem[Aranceta-Bartrina(1999a)]{ref-journal}
% Aranceta-Bartrina, Javier. 1999a. Title of the cited article. \textit{Journal Title} 6: 100--10.
% % Reference 2
% \bibitem[Aranceta-Bartrina(1999b)]{ref-book1}
% Aranceta-Bartrina, Javier. 1999b. Title of the chapter. In \textit{Book Title}, 2nd ed. Edited by Editor 1 and Editor 2. Publication place: Publisher, vol. 3, pp. 54–96.
% % Reference 3
% \bibitem[Baranwal and Munteanu {[1921]}(1955)]{ref-book2}
% Baranwal, Ajay K., and Costea Munteanu. 1955. \textit{Book Title}. Publication place: Publisher, pp. 154--96. First published 1921 (op-tional).
% % Reference 4
% \bibitem[Berry and Smith(1999)]{ref-thesis}
% Berry, Evan, and Amy M. Smith. 1999. Title of Thesis. Level of Thesis, Degree-Granting University, City, Country. Identifi-cation information (if available).
% % Reference 5
% \bibitem[Cojocaru et al.(1999)]{ref-unpublish}
% Cojocaru, Ludmila, Dragos Constatin Sanda, and Eun Kyeong Yun. 1999. Title of Unpublished Work. \textit{Journal Title}, phrase indicating stage of publication.
% % Reference 6
% \bibitem[Driver et al.(2000)]{ref-proceeding}
% Driver, John P., Steffen Rohrs, and Sean Meighoo. 2000. Title of Presentation. In \textit{Title of the Collected Work} (if available). Paper presented at Name of the Conference, Location of Conference, Date of Conference.
% % Reference 7
% \bibitem[Harwood(2008)]{ref-url}
% Harwood, John. 2008. Title of the cited article. Available online: URL (accessed on Day Month Year).
% \end{thebibliography}
% }{}

% % APA format (Used for journal: admsci, behavsci, businesses, econometrics, economies, education, ejihpe, games, humans, ijfs, journalmedia, jrfm, languages, psycholint, publications, tourismhosp, youth)
% \isAPAStyle{%
% \begin{thebibliography}{999}
% % Reference 1
% \bibitem[\protect\citeauthoryear{Azikiwe \BBA\ Bello}{{2020a}}]{ref-journal}
% Azikiwe, H., \& Bello, A. (2020a). Title of the cited article. \textit{Journal Title}, \textit{Volume}(Issue), 
% Firstpage--Lastpage/Article Number.
% % Reference 2
% \bibitem[\protect\citeauthoryear{Azikiwe \BBA\ Bello}{{2020b}}]{ref-book1}
% Azikiwe, H., \& Bello, A. (2020b). \textit{Book title}. Publisher Name.
% % Reference 3
% \bibitem[Davison(1623/2019)]{ref-book2}
% Davison, T. E. (2019). Title of the book chapter. In A. A. Editor (Ed.), \textit{Title of the book: Subtitle} 
% (pp. Firstpage--Lastpage). Publisher Name. (Original work published 1623) (Optional).
% % Reference 4
% \bibitem[Fistek et al.(2017)]{ref-proceeding}
% Fistek, A., Jester, E., \& Sonnenberg, K. (2017, Month Day). Title of contribution [Type of contribution]. Conference Name, Conference City, Conference Country.
% % Reference 5
% \bibitem[Hutcheson(2012)]{ref-thesis}
% Hutcheson, V. H. (2012). \textit{Title of the thesis} [XX Thesis, Name of Institution Awarding the Degree].
% % Reference 6
% \bibitem[Lippincott \& Poindexter(2019)]{ref-unpublish}
% Lippincott, T., \& Poindexter, E. K. (2019). \textit{Title of the unpublished manuscript} [Unpublished manuscript/Manuscript in prepara-tion/Manuscript submitted for publication]. Department Name, Institution Name.
% % Reference 7
% \bibitem[Harwood(2008)]{ref-url}
% Harwood, J. (2008). \textit{Title of the cited article}. Available online: URL (accessed on Day Month Year).
% \end{thebibliography}
% }{}

% If authors have biography, please use the format below
%\section*{Short Biography of Authors}
%\bio
%{\raisebox{-0.35cm}{\includegraphics[width=3.5cm,height=5.3cm,clip,keepaspectratio]{Definitions/author1.pdf}}}
%{\textbf{Firstname Lastname} Biography of first author}
%
%\bio
%{\raisebox{-0.35cm}{\includegraphics[width=3.5cm,height=5.3cm,clip,keepaspectratio]{Definitions/author2.jpg}}}
%{\textbf{Firstname Lastname} Biography of second author}

% For the MDPI journals use author-date citation, please follow the formatting guidelines on http://www.mdpi.com/authors/references
% To cite two works by the same author: \citeauthor{ref-journal-1a} (\citeyear{ref-journal-1a}, \citeyear{ref-journal-1b}). This produces: Whittaker (1967, 1975)
% To cite two works by the same author with specific pages: \citeauthor{ref-journal-3a} (\citeyear{ref-journal-3a}, p. 328; \citeyear{ref-journal-3b}, p.475). This produces: Wong (1999, p. 328; 2000, p. 475)

%%%%%%%%%%%%%%%%%%%%%%%%%%%%%%%%%%%%%%%%%%
%% for journal Sci
%\reviewreports{\\
%Reviewer 1 comments and authors’ response\\
%Reviewer 2 comments and authors’ response\\
%Reviewer 3 comments and authors’ response
%}
%%%%%%%%%%%%%%%%%%%%%%%%%%%%%%%%%%%%%%%%%%
\PublishersNote{}
%\isPreprints{}{% This command is only used for ``preprints''.
\end{adjustwidth}
%} % If the paper is ``preprints'', please uncomment this parenthesis.
\end{document}

